\section{Resultados}
%=====================================================================

\subsection{Detección de QRS sobre MIT-BIH}
En los 15 registros evaluados (Sección~5), la fusión por consenso de las
dos derivaciones alcanzó $Se=P^{+}=F_1\approx 1{,}000$ frente a las
anotaciones de referencia, con un error temporal medio inferior a una
muestra a 360\,Hz. El número de latidos emparejados por registro varió
entre 405 (registro 217, con abundantes latidos de la clase V) y 3\,355
(registro 215), cubriendo un total de 39\,773 latidos emparejados en el proceso. Los detalles por registro están en el informe de rendimiento
individual de cada prueba.

\subsection{Clasificación de latidos (15 pruebas preliminares)}
La Tabla~\ref{tab:campana} resume el desempeño del bosque aleatorio
sobre los 15 registros de MIT-BIH. El \emph{accuracy} global osciló
entre 0,789 (registro 214) y 0,997 (registro 212), con una mediana de
0,945. La clase N se predijo con precisión y \emph{recall} superiores a
0,95 en la mayoría de los registros; la clase V se detectó con
\emph{recall} alto (0,85--1,00) cuando había suficientes ejemplos de
entrenamiento, pero su precisión cayó en registros con morfologías
atípicas (p.\,ej., 214: precisión 0,197 con 517 predicciones V frente a
138 reales). Las clases S y F son las más difíciles: S tiene soportes
pequeños (0--383 latidos) y el bosque tiende a predecir N; F solo
aparece en el registro 208 (353 latidos), donde el \emph{recall} fue
0,448. La \emph{balanced accuracy} (media de \emph{recall} por clase)
osciló entre 0,29 (registro 210) y 0,997 (registro 212); esta métrica
refleja mejor que el \emph{accuracy} el desbalance severo entre clases
(N y V concentran $>90\ porciento$ de los latidos en la mayoría de los
registros).

\begin{table*}[t]
\centering\footnotesize
\caption{Procedimiento de 15 pruebas preliminares sobre MIT-BIH (2
derivaciones, 360\,Hz, bosque de 200 árboles). Se muestran los soportes
por clase en la referencia, las métricas globales y el $F_1$ por clase.
Los registros con soporte 0 en una clase no permiten calcular su
$F_1$ (se indica con «--»).}
\label{tab:campana}
\begin{tabular}{@{}lrrrrrrrrrrrr@{}}
\toprule
 & \multicolumn{4}{c}{Soporte} & \multicolumn{3}{c}{Global} & \multicolumn{4}{c}{$F_1$ por clase}\\
\cmidrule(lr){2-5}\cmidrule(lr){6-8}\cmidrule(lr){9-12}
Reg. & N & S & V & F & Acc. & $F_1$ macro & Bal. acc. & N & S & V & F\\
\midrule
200 & 1742 & 30 & 794 & 2 & 0,942 & 0,470 & 0,463 & 0,959 & 0,000 & 0,921 & 0,000\\
201 & 1635 & 105 & 18 & 0 & 0,945 & 0,529 & 0,680 & 0,971 & 0,174 & 0,971 & --\\
202 & 2060 & 55 & 11 & 1 & 0,970 & 0,380 & 0,341 & 0,985 & 0,000 & 0,533 & 0,000\\
203 & 2491 & 2 & 409 & 1 & 0,975 & 0,473 & 0,458 & 0,986 & 0,000 & 0,908 & 0,000\\
205 & 2570 & 3 & 71 & 11 & 0,992 & 0,486 & 0,475 & 0,996 & 0,000 & 0,948 & 0,000\\
207 & 1540 & 107 & 181 & 0 & 0,973 & 0,694 & 0,875 & 0,984 & 0,856 & 0,935 & --\\
208 & 1579 & 2 & 944 & 353 & 0,931 & 0,639 & 0,612 & 0,944 & 0,000 & 0,994 & 0,617\\
209 & 2620 & 383 & 1 & 0 & 0,895 & 0,311 & 0,393 & 0,943 & 0,302 & 0,000 & --\\
210 & 2421 & 20 & 179 & 9 & 0,932 & 0,311 & 0,291 & 0,964 & 0,000 & 0,279 & 0,000\\
212 & 2746 & 0 & 0 & 0 & 0,997 & 0,250 & 0,997 & 0,998 & -- & -- & --\\
213 & 2640 & 28 & 158 & 361 & 0,877 & 0,403 & 0,497 & 0,959 & 0,000 & 0,653 & 0,000\\
214 & 2001 & 0 & 138 & 1 & 0,789 & 0,297 & 0,511 & 0,875 & -- & 0,311 & 0,000\\
215 & 3188 & 3 & 163 & 1 & 0,968 & 0,377 & 0,339 & 0,983 & 0,000 & 0,525 & 0,000\\
217 & 244 & 0 & 161 & 0 & 0,844 & 0,411 & 0,804 & 0,886 & -- & 0,757 & --\\
219 & 2082 & 7 & 58 & 1 & 0,991 & 0,473 & 0,453 & 0,995 & 0,000 & 0,895 & 0,000\\
\bottomrule
\end{tabular}
\end{table*}

\begin{table*}[t]
\centering\footnotesize
\caption{Tiempos medios (ms) y heap máximo (MiB) por operación,
promediados sobre las 15 pruebas del muestreo (Windows~11, 20
procesadores, OpenJDK~25.0.4.1).}
\label{tab:tiempos}
\begin{tabular}{@{}lrrrr@{}}
\toprule
Operación & Media (ms) & Mín. (ms) & Máx. (ms) & Heap máx. (MiB)\\
\midrule
Leer cabecera & 0,41 & 0,29 & 0,54 & \emph{ver nota}\\
Decodificar señales & 35,6 & 25,1 & 59,9 & \emph{ver nota}\\
Leer anotaciones & 2,06 & 0,81 & 5,15 & \emph{ver nota}\\
Cargar registro & 45,3 & 29,2 & 77,4 & \emph{ver nota}\\
Analizar QRS & 51,6 & 34,2 & 140,9 & \emph{ver nota}\\
Clasificar latidos (RF) & 144,9 & 97,3 & 204,7 & \emph{ver nota}\\
Abrir visor & 30,4 & 22,3 & 135,9 & \emph{ver nota}\\
Exportar PNG & 37,2 & 32,9 & 67,3 & \emph{ver nota}\\
\bottomrule
\end{tabular}
\end{table*}

\emph{Nota:} el heap máximo de la Tabla~\ref{tab:tiempos} se informa en
el Excel por prueba y no por operación; los valores oscilaron entre
20,6\,MiB (P001, lectura de cabecera) y 837,7\,MiB (P006, análisis QRS),
con una tendencia creciente a lo largo de la sesión por acumulación de
referencias vivas entre pruebas (la JVM no se reinició entre P001 y
P015). El detalle por prueba está en el archivo
\texttt{resultados\_pruebas.xlsx} y en los informes PDF individuales.



\subsection{Ejemplo de informe de rendimiento}
La evaluación experimental de un registro de 60\,s y 12 derivaciones
(tres repeticiones tras el calentamiento) produjo un informe PDF de tres
páginas. Las operaciones de mayor coste fueron la exportación PNG de la
imagen estándar de $1\,891\times852$ píxeles (125\,ms de media) y su
renderizado (105\,ms), frente a 39\,ms para cargar el registro (de ellos
25\,ms de decodificación) y 28\,ms para la detección en las doce
derivaciones. En ese informe el heap máximo observado fue de 29,3\,MiB
mientras que la memoria residente fue de 103\,MiB, una diferencia de más
de tres veces que ilustra por qué ambas magnitudes se informan por
separado. En la plataforma Linux utilizada fue posible informar CPU de
usuario y de sistema; en Windows las columnas de CPU aparecen como «no
disponible» y solo se informa el heap (como en la
Tabla~\ref{tab:tiempos}).
